% !TEX program = pdflatex
% Twelve logical slides following presentation/outline.md, including the title slide.
% The embedding quiz has two overlays, producing thirteen PDF pages.
% Build from the repository root: ./scripts/build-tex.sh presentation/local-ai-workshop.tex
% latexmk resolves the total slide count. The earlier full deck is retained on the main branch.
\documentclass[aspectratio=169,12pt]{beamer}
\usetheme{Madrid}
\usecolortheme{default}
\setbeamertemplate{navigation symbols}{}
\setbeamertemplate{footline}[frame number]
\setbeamertemplate{itemize items}[circle]
\setbeameroption{hide notes}
\usepackage[T1]{fontenc}
\usepackage{lmodern}
\usepackage{booktabs}
\usepackage{array}
\usepackage{tabularx}
\usepackage{tikz}
\usetikzlibrary{arrows.meta,positioning,calc}
\newcolumntype{Y}{>{\raggedright\arraybackslash}X}
\newcommand{\source}[1]{\par\vfill{\fontsize{7}{8.5}\selectfont\color{black!65}#1\par}}
\newcommand{\modelid}[1]{{\fontsize{8}{9.5}\selectfont\texttt{#1}}}
\newcommand{\lead}[1]{{\usebeamercolor[fg]{structure}\bfseries #1}}
\title{DIY-AI}
\subtitle{Running Large Language and Agentic Models\\in your own environments}
\author{Dr. Coert van Gemeren}
\institute{HU University of Applied Sciences Utrecht}
\date{26 October 2026}
\hypersetup{pdftitle={DIY-AI: Running Large Language and Agentic Models in your own environments},
  pdfauthor={Dr. Coert van Gemeren},
  pdfsubject={Twelve-slide workshop with an embedding-space quiz based on outline.md}}

\begin{document}

% SLIDE 1 -------------------------------------------------------------
\begin{frame}
\titlepage
\note{Title, subtitle, presenter, institution and event date follow
presentation/outline.md slide 1.}
\end{frame}

% SLIDE 2 -------------------------------------------------------------
\begin{frame}{Local AI: models, runtimes and hardware}
\small
A trained \textbf{model} supplies the weights. A \textbf{runtime} loads them
and runs inference on compatible \textbf{hardware}.
\par\medskip

{\footnotesize
\begin{tabularx}{\textwidth}{@{}>{\raggedright\arraybackslash}p{0.19\textwidth}Y >{\raggedright\arraybackslash}p{0.22\textwidth}@{}}
\toprule
\textbf{Runtime} & \textbf{Example model} & \textbf{Workshop use} \\
\midrule
Ollama service & \modelid{gemma4:26b} & Text chat on Mac \\
\addlinespace[0.3em]
vLLM Metal service & \modelid{mlx-community/gemma-4-e4b-it-4bit} & Text, no audio in this setup \\
\addlinespace[0.3em]
vLLM Metal service & \modelid{openai/whisper-large-v3} & Audio transcription \\
\addlinespace[0.3em]
MLX-VLM program & \modelid{google/gemma-4-E4B-it} & Audio input on Mac \\
\addlinespace[0.3em]
vLLM service & \modelid{deepseek-ai/DeepSeek-V4-Flash} & Text and code on Spark \\
\bottomrule
\end{tabularx}}
\par\medskip
\lead{Local AI is practical:} match your task and model to the runtime and available memory.
\source{Workshop configurations from the outline, notebooks and launch scripts.
Runtime references: \href{https://github.com/vllm-project/vllm-metal}{vLLM Metal},
\href{https://github.com/Blaizzy/mlx-vlm}{MLX-VLM},
\href{https://docs.ollama.com/quickstart}{Ollama}.}
\note{Sources: presentation/outline.md slide 2; scripts/start-mlx-vllm.sh; scripts/start-mlx-vlm.sh;
notebooks/local-llm-vllm.ipynb, notebooks/local-llm-ollama.ipynb,
notebooks/local-llm-whisper.ipynb, notebooks/local-vlm-gemma.ipynb,
notebooks/local-llm-spark-cluster.ipynb.
The no-audio label describes the workshop vLLM Gemma configuration, not every
Gemma 4 model or current vLLM implementation. vLLM Metal is the Apple Silicon
plugin using MLX. References checked 21 October 2026:
\url{https://github.com/vllm-project/vllm-metal};
\url{https://github.com/Blaizzy/mlx-vlm}; \url{https://docs.ollama.com/quickstart}.}
\end{frame}

% SLIDE 3 -------------------------------------------------------------
\begin{frame}{Four ways to run the workshop models}
\small
\begin{tabularx}{\textwidth}{@{}>{\raggedright\arraybackslash}p{0.27\textwidth}Y >{\raggedright\arraybackslash}p{0.27\textwidth}@{}}
\toprule
\textbf{Hardware and runtime} & \textbf{How we use it} & \textbf{Entry point} \\
\midrule
\textbf{Apple Silicon + Ollama} & Easiest starting point for local chat & \modelid{ollama run gemma4:26b} \\
\addlinespace[0.3em]
\textbf{Apple Silicon + vLLM Metal} & A local model service for notebooks and apps & \modelid{./scripts/start-mlx-vllm.sh}\newline API at \modelid{localhost:8000} \\
\addlinespace[0.3em]
\textbf{Apple Silicon + MLX-VLM} & Direct Python inference, including Gemma audio & \modelid{notebooks/}\newline\modelid{local-vlm-gemma.ipynb} \\
\addlinespace[0.3em]
\textbf{DGX Spark + vLLM} & Connect over VPN and call the remote model service & Change the API base URL to Spark \\
\bottomrule
\end{tabularx}
\par\medskip
After installation and download, start a model with a command or a cell.\par
Gemma's saved local model load: \textbf{4.8 seconds}.
\source{Sources: local launch scripts and notebooks. Spark audio uses 1 DGX Spark.
DeepSeek uses 2 DGX Sparks. Hardware and VPN configuration follow the workshop outline.}
\note{Sources: presentation/outline.md slide 3; scripts/start-mlx-vllm.sh;
notebooks/local-vlm-gemma.ipynb cell index 2 records a 4.8 s load;
notebooks/local-llm-qwen-omni.ipynb uses spark-0240:8000;
benchmark-results/spark.json uses spark-0240:9000.
First-time installation, downloads, service startup and model loading are distinct
from warm inference. These launch commands assume setup is complete.}
\end{frame}

% SLIDE 4 -------------------------------------------------------------
\begin{frame}{Three tests of useful model behavior}
\small
\lead{1. Five questions: responsiveness and completeness}\par
Capital of France, an AI poem, quantum computing, the laws of robotics
and the plot of \emph{Inception}.\par
\textbf{Measures:} request latency, completion tokens per second and cut-off answers.
\par\bigskip

\lead{2. Dutch transcription: agreement with a curated reference}\par
Jeroen Hendrikse's radiology recording from UMC Utrecht.\par
\textbf{Measure:} word error rate (WER) over 1,350 reference words.
\par\bigskip

\lead{3. LiveCodeBench-Pro: executable programming solutions}\par
Ten contest problems, with one C++ solution attempt per problem.\par
\textbf{Measure:} accepted solutions out of ten selected problems.
\par\medskip

\textbf{Only DeepSeek on Spark produced judgeable code among our self-hosted setups.}
GPT-6 Astra provides a separate cloud comparison.
\source{Sources: original presentation, transcription notebooks and saved LiveCodeBench-Pro results.
The five-question test is a smoke test, without a systematic correctness score.}
\note{Five-question prompts and metrics reuse the historical full deck,
available as presentation/presentation.tex on the main branch.
Audio: notebooks/local-llm-whisper.ipynb, notebooks/local-vlm-gemma.ipynb,
notebooks/local-llm-qwen-omni.ipynb.
Coding: benchmark-results/ollama.json and benchmark-results/spark.json.
Ollama has 10 No Code responses. Local vLLM was not run in the saved judged comparison.
GPT-6 Astra's later Codex-session evaluation is recorded in
presentation/livecodebench-gpt-6-astra-results.tex and
benchmark-results/gpt-6-astra-2026-10-21/. Its generation conditions differ.}
\end{frame}

% SLIDE 5 -------------------------------------------------------------
\begin{frame}{Five questions: speed and completion}
\small
\centering
\begin{tabular}{@{}lrrr@{}}
\toprule
\textbf{Request latency} & \textbf{Local vLLM} & \textbf{Ollama} & \textbf{Spark} \\
 & Gemma E4B & Gemma 26B & DeepSeek V4 Flash \\
\midrule
Capital of France & 1.745\,s & 1.478\,s & \textbf{0.607\,s} \\
AI poem & 5.491\,s & 19.619\,s & \textbf{3.686\,s} \\
Quantum computing & 7.818\,s & 22.523\,s & 23.756\,s \\
Laws of robotics & 7.463\,s & 14.958\,s & 8.349\,s \\
\emph{Inception} & 7.794\,s & 24.565\,s & 15.920\,s \\
\midrule
Mean latency & 6.062\,s & 16.628\,s & 10.464\,s \\
Completion rate (tokens/s) & 18.64 & 34.37 & 28.94 \\
Incomplete answers & 3 of 5 & 2 of 5 & 0 of 5 \\
\bottomrule
\end{tabular}
\par\medskip
\raggedright
Local vLLM's three faster answers were cut short.\par
Spark finished all five responses with \texttt{stop}.
\source{Reused from the original deck's 7 October 2026 snapshot.
Token budgets: 150 / 800 / 800. Different models, tokenizers and reasoning limit comparisons.
Completion rate = total completion tokens / total request time.}
\note{Source: historical presentation/presentation.tex on the main branch, frames Per-Request End-to-End Latency,
Latency Summary, Token Counts and Throughput, and Response Completeness Matters.
The user specifically requested the original presentation's first-test results.
The current notebooks/local-llm-spark-cluster.ipynb contains a newer five-question run,
so this slide deliberately retains the dated original snapshot.
The historical local vLLM notebook on the main branch did not retain finish
reasons. Three visible answers end mid-response at the 150-token limit.
The current notebook collects new measurements independently of Spark. Ollama has two length finish reasons.
Spark has five stop finish reasons. Completion does not establish correctness.}
\end{frame}

% SLIDE 6 -------------------------------------------------------------
\begin{frame}{Dutch transcription: three models on one recording}
\small
Jeroen Hendrikse, UMC Utrecht: \emph{radiologieverslag}\par
\textbf{473.6 seconds of speech}, scored against 1,350 curated reference words.
\par\medskip

\begin{tabularx}{\textwidth}{@{}Y p{0.23\textwidth}rr@{}}
\toprule
\textbf{Model} & \textbf{Runtime / hardware} & \textbf{WER} $\downarrow$ & \textbf{Time} \\
\midrule
Whisper large-v3\newline\modelid{openai/whisper-large-v3}
 & vLLM Metal\newline Apple Silicon & 12.00\% & Not recorded \\
\addlinespace[0.8em]
\textbf{Gemma 4 E4B}\newline\modelid{google/gemma-4-E4B-it}
 & MLX-VLM\newline Apple Silicon & \textbf{7.78\%} & 144.3\,s \\
\addlinespace[0.8em]
Qwen3-Omni 30B A3B\newline\modelid{Qwen/Qwen3-Omni-30B-A3B-Instruct}
 & vLLM\newline 1 DGX Spark & 8.44\% & 98.1\,s \\
\bottomrule
\end{tabularx}
\par\medskip
Gemma has the lowest WER. Qwen finishes the transcription loop sooner.
\source{Sources: saved outputs of notebooks/local-llm-whisper.ipynb,
notebooks/local-vlm-gemma.ipynb and notebooks/local-llm-qwen-omni.ipynb.
Time covers the transcription loop, excluding model loading and
initial audio decoding. Whisper's notebook does not time the request.}
\note{All scores verified against current saved outputs on 21 October 2026.
Whisper cell index 7: 162/1350 = 12.00 percent.
Gemma cell index 7: 105/1350 = 7.78 percent; cell 3: 144.3 seconds.
Qwen cell index 7: 114/1350 = 8.44 percent; cell 3: 98.1 seconds.
Gemma and Qwen use the same 18 chunks, maximum 28 seconds, and the same prompt.
Whisper submits the complete MP3 to /v1/audio/transcriptions with language nl.
All use temperature 0. These are measured pipelines with different hardware and
preprocessing paths, not an isolated test of model size or backend speed.}
\end{frame}

% SLIDE 7 -------------------------------------------------------------
\begin{frame}{Word error rate: Gemma's small advantage}
\small
\[
\mathrm{WER}=\frac{\text{substitutions}+\text{deletions}+\text{insertions}}
{\text{reference words}}\times100\%
\]
\centering
\begin{tabular}{@{}lrrrr@{}}
\toprule
\textbf{Model} & \textbf{Substitutions} & \textbf{Deletions} & \textbf{Insertions} & \textbf{Total} \\
\midrule
Whisper large-v3 & 86 & 41 & 35 & 162 \\
\textbf{Gemma 4 E4B} & 69 & 21 & 15 & \textbf{105} \\
Qwen3-Omni & 73 & 15 & 26 & 114 \\
\bottomrule
\end{tabular}
\par\bigskip
\raggedright
\lead{Gemma: 9 fewer word errors than Qwen}\par
7.78\% versus 8.44\% WER: a \textbf{0.67 percentage-point} gap
using the unrounded scores.
\par\medskip

Scoring ignores case and punctuation, splits hyphens, and retains fillers
and repetitions. Curated TXT and SRT references yield identical scores.
\source{Sources: the three transcription notebooks, cell index 7. One recording and different
inference pipelines. This small difference does not establish a general model ranking.
WER counts word differences without weighting their clinical meaning.}
\note{Reuses and updates the original Audio Evaluation slide.
Absolute gap: (114-105)/1350*100 = 0.6666667 percentage points.
Relative error reduction: 9/114 = 7.8947 percent (not displayed).
The displayed WER values round independently, so their difference rounds to 0.66;
the slide's 0.67 pp uses the underlying integer counts.
The scoring routine implements word-level minimum edit distance after NFKC
normalization and case folding. Reference editorial choices can affect WER.}
\end{frame}

% SLIDE 8 -------------------------------------------------------------
\begin{frame}{LiveCodeBench-Pro: a harder programming test}
\small
Contest problems from \textbf{Codeforces, ICPC and IOI} require an algorithm
and working code that passes the judge's tests.
\par\medskip

\begin{columns}[T,onlytextwidth]
\begin{column}{0.51\textwidth}
\lead{What happened locally}\par\smallskip
Local MLX trials did not yield usable solutions in the workshop.\par
The saved comparison records:
\begin{itemize}
\setlength{\itemsep}{0.4em}
\item \textbf{Ollama, Gemma 26B:} all 10 responses hit the 4,096-token limit
without final code.
\item \textbf{Local vLLM: Not run.}
\end{itemize}
\end{column}
\begin{column}{0.44\textwidth}
\lead{The larger model}\par\smallskip
\textbf{DeepSeek V4 Flash}\par
\modelid{deepseek-ai/DeepSeek-V4-Flash}
\par\medskip
vLLM over VPN on \textbf{2 DGX Sparks}\par
\textbf{256 GB} total unified memory\par
(128 GB per node)
\par\medskip
One C++ attempt per problem.\par
Temperature 0, 4,096 tokens.
\end{column}
\end{columns}
\par\medskip
The trials expose task difficulty and a token-budget limit, without
establishing a general limit for local models.
\source{Local MLX observation and cluster configuration: workshop outline.
Saved evidence: Ollama, 8 October; Spark, 20 October 2026.
\href{https://arxiv.org/abs/2506.11928}{LiveCodeBench-Pro paper};
\href{https://www.nvidia.com/en-us/products/workstations/dgx-spark/}{DGX Spark specifications}.}
\note{The user reports unsuccessful local MLX trials in outline.md slide 8.
The retained benchmark artifacts do not contain a completed local MLX/vLLM run.
Do not convert Not run to a measured zero or attribute Ollama's No Code outcomes to MLX.
benchmark-results/ollama.json has 10 No Code verdicts and 10 length finish reasons.
The benchmark uses the first 10 unique problem IDs in sorted order, seven easy
and three medium; it contains no hard-tier problems. Native ARM64 judging.
Dataset revision adebffce047dddb7768a86bace6aea4f7425e3bc.
Upstream commit d614eef4f95b3654bf0758c64e1194b2575850b8.
Sources checked 21 October 2026: \url{https://arxiv.org/abs/2506.11928};
\url{https://www.nvidia.com/en-us/products/workstations/dgx-spark/}.
256 GB is the aggregate across two 128 GB nodes, not a single shared address space.}
\end{frame}

% SLIDE 9 -------------------------------------------------------------
\begin{frame}{LiveCodeBench-Pro: Spark and GPT-6 Astra}
\small
Same ten problems and ARM64 judge. Both runs completed judging for all ten.
\par\medskip
\centering
\begin{tabular}{@{}lcc@{}}
\toprule
 & \textbf{DeepSeek V4 Flash} & \textbf{GPT-6 Astra} \\
 & 2 DGX Sparks, 20 October & Codex session, 21 October \\
\midrule
Accepted / selected & \textbf{4/10 (40\%)} & \textbf{10/10 (100\%)} \\
Compilation error & 4 & 0 \\
Wrong answer & 2 & 0 \\
Mean inference time & 25.32\,s & Not measured \\
\bottomrule
\end{tabular}
\par\medskip
\raggedright
\lead{Astra passed all ten problems.}
DeepSeek's six failures included four compilation errors.
\par\medskip
\textbf{Generation conditions differ:}\par
DeepSeek: temperature 0, 4,096-token API budget per problem.\par
Astra: Codex with \texttt{xhigh} reasoning, shared conversation context,
no per-problem token cap.
\source{Source: livecodebench-gpt-6-astra-results.tex and saved JSON.
One judged submission per problem. Astra sources frozen before testing, with no revisions
after feedback. Subset: 7 easy, 3 medium, 0 hard. Generation conditions differ.
These are workshop results, not full-benchmark leaderboard scores.}
\note{Source: presentation/livecodebench-gpt-6-astra-results.tex, corroborated on
22 October 2026 against benchmark-results/gpt-6-astra-2026-10-21/gpt-6-astra.json,
deepseek-baseline.json and manifest.json. Both models use the same ten problem IDs,
dataset revision adebffce047dddb7768a86bace6aea4f7425e3bc, upstream commit
d614eef4f95b3654bf0758c64e1194b2575850b8, and linux/arm64 LightCPVerifier judge.
DeepSeek: 20 October 2026, 4 Accepted, 4 Compilation Error, 2 Wrong Answer.
Mean inference latency 25.3208730209 s, no truncated responses.
GPT-6 Astra: 21 October 2026, ten Accepted verdicts and 582/582 recorded judge
cases passed. Sources were authored and reviewed in a shared Codex conversation
with xhigh reasoning, frozen before any sample execution or judge submission,
and submitted once each. No revisions after feedback, no editorials, reference
solutions or hidden testcase inputs consulted. No per-problem token cap,
temperature setting, inference latency or token-usage measurements.
The Codex-session result is not a controlled API comparison with DeepSeek.
It replaces the earlier historical cloud-model context with the user's own
evaluation of the same subset. Model IDs: gpt-6-astra and
deepseek-ai/DeepSeek-V4-Flash.}
\end{frame}

% SLIDE 10 -------------------------------------------------------------
\begin{frame}{Local document search with RAG on DGX Spark}
\small
Retrieval-augmented generation (RAG) supplies relevant document passages
to the model before it writes an answer.
\par\medskip

% Editable TikZ diagram explicitly requested by outline.md.
\centering
\begin{tikzpicture}[
  x=1cm,y=1cm,
  ragnode/.style={draw=structure.fg,fill=white,align=center,
    minimum height=0.65cm,text width=2.4cm,font=\small,inner sep=4pt},
  flow/.style={-{Stealth[length=2mm]},thick,draw=structure.fg},
  raglabel/.style={font=\scriptsize,align=center,text=black!70}]
\node[ragnode] (docs) at (0,0) {Documents};
\node[ragnode] (chunks) at (3.2,0) {Chunking};
\node[ragnode,text width=2.6cm,font=\fontsize{9}{10}\selectfont] (embed) at (6.4,0) {EmbeddingGemma};
\node[ragnode] (db) at (9.6,0) {Vector DB\,(Milvus)};
\draw[flow] (docs) -- (chunks);
\draw[flow] (chunks) -- (embed);
\draw[flow] (embed) -- (db);

\node[ragnode] (query) at (0,-1.3) {Question};
\node[ragnode,text width=2.6cm,font=\fontsize{9}{10}\selectfont] (qembed) at (3.2,-1.3) {EmbeddingGemma};
\draw[flow] (query) -- (qembed);
\draw[flow] (qembed.east) -- node[raglabel,above,pos=0.55] {query vector} (9.6,-1.3) -- (db.south);

\node[ragnode,dashed,text width=3.65cm] (rerank) at (8.9,-2.65)
  {BGE reranker v2.5\\Gemma2 lightweight\\[-1pt]{\scriptsize optional}};
\node[ragnode] (llm) at (4.4,-2.65) {Gemma 4 31B};
\node[ragnode] (answer) at (0,-2.65) {Answer\\with source passages};
\draw[flow] (db.east) -- (11.15,0) --
  node[raglabel,right] {Top-K} (11.15,-2.65) -- (rerank.east);
\draw[flow] (rerank.west) -- node[raglabel,above] {Top-N} (llm.east);
\draw[flow] (llm.west) -- (answer.east);
\draw[flow] (query.west) -- (-1.6,-1.3) -- (-1.6,-3.5) --
  node[raglabel,below,pos=0.55] {original question} (4.4,-3.5) -- (llm.south);
\end{tikzpicture}
\par\smallskip
\raggedright
Top-K retrieves candidate passages. The optional reranker selects Top-N.
Without reranking, the retrieved passages go directly to Gemma.
\source{Proposed local stack from the outline, not a benchmarked deployment.
Model references: \href{https://huggingface.co/google/embeddinggemma-300m}{EmbeddingGemma 300M},
\href{https://huggingface.co/BAAI/bge-reranker-v2.5-gemma2-lightweight}{BGE reranker},
\href{https://huggingface.co/google/gemma-4-31B-it}{Gemma 4 31B};
\href{https://milvus.io/docs/quickstart.md}{Milvus documentation}.}
\note{Architecture follows outline.md slide 10. The query also needs an embedding
in the same vector space as document chunks. Retain document text and source
metadata with vectors. The reranker scores query/passage pairs, then the generator
receives the original question and selected passages. Diagram arrow Top-K carries
passages and their metadata, not just vectors. The original question also goes
to the reranker (omitted as a separate connector for legibility).
Model IDs: google/embeddinggemma-300m;
BAAI/bge-reranker-v2.5-gemma2-lightweight; google/gemma-4-31B-it.
The BGE card identifies this model as a Gemma2-9B-based reranker with layer and
token compression. The diagram is a proposed architecture, without claims of
measured retrieval quality, simultaneous memory fit, latency or clinical validation.
Sources checked 21 October 2026:
\url{https://huggingface.co/google/embeddinggemma-300m};
\url{https://huggingface.co/BAAI/bge-reranker-v2.5-gemma2-lightweight};
\url{https://huggingface.co/google/gemma-4-31B-it};
\url{https://milvus.io/docs/quickstart.md}.}
\end{frame}

% SLIDE 11 -------------------------------------------------------------
\begin{frame}{Embeddings: meaning as coordinates}
\small
\begin{columns}[T,onlytextwidth]
\begin{column}{0.36\textwidth}
\lead{Text becomes a vector}\par\smallskip
An embedding is a \textbf{list of numbers} representing text.
\par\smallskip
Similar meanings tend to lie nearby in a high-dimensional space.
\par\smallskip

\lead{Meaning has direction}\par\smallskip
{\footnotesize
$E(\text{King}) + E(\text{Woman}) - E(\text{Man})$\par
\hspace*{0.8cm}$\approx E(\text{Queen})$\par}\smallskip
{\scriptsize $E(\text{word})$ is its embedding vector.\par}
\end{column}
\begin{column}{0.61\textwidth}
\centering
\lead{Quiz: which country?}\par\smallskip
{\normalsize
$\text{Holland} + \text{Fries} - \text{Kroket} \approx {}$
\makebox[1.55cm][l]{\strut\alt<2>{\textbf{Belgium}}{\textbf{?}}}\par}
\smallskip

% Invented coordinates for an explanatory diagram, not measured embeddings.
% Fries - Kroket = (4.3, 0.45) = result - Holland.
% The fixed bounding box and fixed-width answer keep both overlays aligned.
\begin{tikzpicture}[
  x=1.08cm,y=0.90cm,
  vector/.style={-{Stealth[length=1.5mm]},thin,draw=black!25},
  shiftvector/.style={-{Stealth[length=2.2mm]},very thick,draw=structure.fg},
  point/.style={circle,fill=structure.fg,inner sep=2pt},
  word/.style={font=\small,text=structure.fg},
  shiftlabel/.style={font=\scriptsize,fill=white,inner sep=2pt,text=structure.fg}]
\path[use as bounding box] (-0.2,-0.35) rectangle (7.0,3.4);
\draw[step=0.5,black!7,very thin] (0,0) grid (6.75,3.25);
\draw[-{Stealth[length=1.4mm]},black!30] (0,0) -- (6.9,0);
\draw[-{Stealth[length=1.4mm]},black!30] (0,0) -- (0,3.35);
\node[font=\tiny,text=black!50,below left] at (0,0) {$0$};
\coordinate (kroket) at (1.0,0.65);
\coordinate (fries) at (5.3,1.1);
\coordinate (holland) at (1.6,2.25);
\coordinate (result) at (5.9,2.7);
\foreach \wordpoint in {kroket,fries,holland,result}
  \draw[vector] (0,0) -- (\wordpoint);
\draw[shiftvector] (kroket) --
  node[shiftlabel,above,sloped] {$+\,\text{Fries}\;-\;\text{Kroket}$} (fries);
\draw[shiftvector] (holland) --
  node[shiftlabel,above,sloped] {$+\,\text{Fries}\;-\;\text{Kroket}$} (result);
\node[point] at (kroket) {};
\node[point] at (fries) {};
\node[point] at (holland) {};
\node[point] at (result) {};
\node[word,below=3pt] at (kroket) {Kroket};
\node[word,below=3pt] at (fries) {Fries};
\node[word,above=4pt] at (holland) {Holland};
\node[word,above=4pt,text width=1.55cm,align=center,font=\small\bfseries]
  at (result) {\strut\alt<2>{Belgium}{?}};
\end{tikzpicture}
\par
{\fontsize{8}{9.5}\selectfont\color{black!65}
Illustrative 2D embedding space; not a measured model result.\par}
\end{column}
\end{columns}
\par\smallskip

\begin{overlayarea}{\textwidth}{0.55cm}
\centering
\only<2>{\textbf{Add the fries association and subtract the kroket association.}}
\end{overlayarea}
\lead{Back to RAG:} EmbeddingGemma represents questions and document chunks
in a shared space so similarity search can retrieve relevant passages.
\source{Embedding intuition and word analogy:
\href{https://www.3blue1brown.com/lessons/gpt/}{3Blue1Brown, Transformers, the tech behind LLMs}.
Text retrieval: \href{https://huggingface.co/google/embeddinggemma-300m}{EmbeddingGemma model card}.}
\note{Slide 11 follows the RAG architecture with an embedding-space explanation.
First explain an embedding as a list of numbers, interpreted as coordinates in
a high-dimensional space. Similar meanings tend to be nearby; some learned
directions express relationships. 3Blue1Brown illustrates this using
E(King) + E(Woman) - E(Man) approximately E(Queen), while explicitly noting
that the analogy is imperfect. E denotes the embedding function. These are word-embedding intuitions, not a guarantee
that all embedding models obey semantic arithmetic.
Ask the audience: Holland + Fries - Kroket is approximately which country?
Pause for answers, then advance once to reveal Belgium. The intended association
replaces kroket with fries while retaining the country role. The coordinates and
snack associations are invented to illustrate the vector operation, not measured
EmbeddingGemma outputs or exclusive cultural associations. Both arrows show
exactly the same displacement in this toy diagram. The real space has many more
than two dimensions, and analogy results need not land exactly on a word vector.
EmbeddingGemma embeds whole questions and document chunks for retrieval, rather
than merely looking up individual token vectors. Query and document encoding
use the model's respective task prompts and produce vectors in a shared space.
Retrieval ranks passages by vector similarity; it does not require snack-style
vector arithmetic. Both overlay pages retain logical slide number 11 of 12.
References checked 23 October 2026:
\url{https://www.3blue1brown.com/lessons/gpt/};
\url{https://huggingface.co/google/embeddinggemma-300m}.}
\end{frame}

% SLIDE 12 -------------------------------------------------------------
\begin{frame}{Next: agentic orchestration in clinical imaging}
\begin{columns}[c,onlytextwidth]
\begin{column}{0.53\textwidth}
{\Large\usebeamercolor[fg]{structure}Local models and document retrieval
are building blocks for a wider workflow.\par}
\par\bigskip
The next presenter will introduce an agentic orchestrator
for use in clinical imaging practice.
\par\bigskip
\small
Coordinating model calls, retrieved information and tools
around the steps of a clinical imaging workflow.
\par\bigskip
{\large\textbf{Over to the next presenter.}}
\end{column}
\begin{column}{0.44\textwidth}
\centering
\small\lead{Try it after the presentation}\par\smallskip
% Generated with scripts/generate-workshop-qr.py; four-module white quiet zone.
\href{https://github.com/uashogeschoolutrecht/IMAGINE-workshop-2026}{%
  \includegraphics[width=6.0cm]{images/github-link.pdf}}
\par\smallskip
\textbf{Scan to download the workshop}\par
\end{column}
\end{columns}
\note{Handover requested in presentation/outline.md slide 12.
No presenter name, product name or deployment claims were supplied.
Invite the audience to download the workshop using the supplied QR code
and try the examples after the presentation. Keep the closing slide visible
long enough for the audience to scan. The QR code is also a clickable link.}
\end{frame}

\end{document}
